% Conclusion section for OrbitGraph paper.

\section{Conclusion}
\label{sec:conclusion}

Low Earth orbit mega-constellations are, at their core, dynamic graphs:
vertices are satellites and ground assets, and edges appear and vanish as
orbital geometry evolves. Routing here is not a one-time shortest-path problem
but a continuous sequence of graph snapshots $G_t$ that must reach the
forwarding plane faster than contacts change.
Distributed link-state protocols such as OSPF were designed for topologies that
change sporadically. They discover mutations reactively and reconverge after the
fact, which our measurements show can black-hole traffic for many seconds at
each ground--satellite handover.

Software-defined networking offers a different contract.
By decoupling the control plane from the data plane, OrbitGraph
centralizes graph construction from ephemeris-driven delay
matrices and multi-source Dijkstra while leaving forwarding to the
existing Linux kernel on each node.
The data plane executes simple, fast operations such as longest-prefix match
and local delivery, while the control plane owns global visibility, policy, and
when updates land relative to link lifetime.
That separation is efficient in two senses. Algorithmically, recomputation stays
in the millisecond range even at 102 nodes. Operationally, only changed
next-hops are programmed, ${\sim}10^2$–$10^3$ routes at handover versus
${\sim}10^4$ at initialization, avoiding the cost of re-flooding a full
link-state database to every router on every orbital tick.

The decisive advantage, however, is not raw compute speed but control
over update ordering.
OrbitGraph's proactive handover installs the post-contact FIB while both old and
new ground--satellite links coexist, make-before-break at the topology level.
Distributed OSPF cannot match that behavior without protocol extensions,
because no single router observes the full mutation schedule in advance.
Our container-based evaluation on StarryNet~\cite{lai2023starrynet}, spanning
27 to 102 nodes over ten repetitions, shows 0\,ms data-plane
outage on every coverage-preserving GSL handover against OSPF black holes
reaching ${\sim}15$\,s, 0\% packet loss through handover and post-handover under
OrbitGraph while OSPF loss reaches 100\% on several grids, comparable
steady-state latency, and incremental updates that keep programmatic cost
proportional to handover scope rather than constellation size. Where the
constellation is too sparse to keep a ground station in view, a coverage gap
severs the path for both protocols, delimiting the regime in which control-plane
ordering is decisive.

OrbitGraph demonstrates that graph-centric SDN routing keeps an emulated LEO
network connected through predictable topology changes, the regime where orbital
mechanics already give operators a global view and a contact
schedule~\cite{kassing2020exploring,handley:ground-relays:2019}.
It is not a universal replacement for distributed routing. Under unplanned bulk
failures with our container-based install path, centralized updates can lag
OSPF, and reaching thousands of nodes will require hierarchical control and
faster southbounds.
Within its effective range of scheduled GSL handovers and orbit-predictable
graph snapshots, OrbitGraph shows that treating the constellation as an
explicit, time-indexed graph and programming forwarding from a centralized
controller is a practical path to the sub-second availability that
latency-sensitive space--terrestrial services require.

From a graph-centric computing view, when the network is a graph that changes
on a known timetable, the effective model makes that graph a first-class control
object: recompute on snapshots, push structured deltas, and align commits with
graph edits, rather than forcing each node to infer global state from local
messages.
OrbitGraph is a concrete instantiation of that idea for LEO routing, validated
end-to-end in a real TCP/IP emulation environment with measurable gains in
data-plane availability where it matters most.
